\chapter{Methodology}

Comparing four extraction strategies fairly requires a controlled setup: the
same documents, the same scoring framework, and a baseline simple enough to
serve as a meaningful lower bound. Figure~\ref{fig:eval_pipeline} shows the
complete evaluation pipeline.

\begin{figure}[h!]
    \centering
    \includegraphics[width=\textwidth,height=0.45\textheight,keepaspectratio]{eval_pipeline.png}
    \caption{Evaluation pipeline: corpus, strategies, ground truth, metrics, and final scores}
    \label{fig:eval_pipeline}
\end{figure}

\section{Research Design}

This thesis uses a comparative study to evaluate four PDF extraction strategies
against a shared test corpus and scoring framework. The comparative approach is
standard for PDF parsing research. Meuschke et al.\ \autocite{meuschke2023}
applied it to benchmark PDF tools across multiple document domains. Ingemarsson
and Persson \autocite{ingemarsson2024} used the same design to rank model families
by accuracy and computational cost. Using a shared methodology means the results
sit alongside that published work rather than in isolation.

The four strategies are text-only, image-only, hybrid, and adaptive. The first
three are fixed: every page goes through the same pathway regardless of content
type. The adaptive system classifies each page first, then routes it to the
appropriate extraction method. This is the only design that can skip the LLM
entirely on plain-text pages. That is what RQ3 tests.

The five research questions each target a different aspect of this comparison.
RQ1 and RQ2 compare accuracy across all four strategies: RQ1 asks whether the
adaptive system beats every fixed strategy, and RQ2 asks how the fixed
strategies rank against each other across document categories. RQ3 covers
efficiency: processing time, token usage, and LLM call count per strategy. RQ4
drives the metric design in Section~\ref{sec:eval_metrics}. RQ5 checks whether the strategy
rankings stay consistent across different thesis documents.

Markdown is the target output format for all strategies. Unlike PDF or Word, it
stores structure in plain text, requires no proprietary software to read, and
works directly with LLM and RAG systems \autocite{cone2019, khan2024}.
Table~\ref{tab:format_comparison} compares all three formats across eight
properties.

\begin{table}[ht]
\centering
\caption{Comparison of PDF, Word (.docx), and Markdown as document formats}
\label{tab:format_comparison}
\renewcommand{\arraystretch}{1.4}
\begin{tabularx}{\textwidth}{lXXX}
\toprule
\textbf{Property} & \textbf{PDF} & \textbf{Word (.docx)} & \textbf{Markdown} \\
\midrule
Human readable (raw) & No & No & \textbf{Yes} \\
Machine parseable & Difficult & Moderate & \textbf{Easy} \\
Version control (Git) & Poor & Poor & \textbf{Excellent} \\
Format portability & High & Low & \textbf{High} \\
Preserves visual fidelity & Excellent & Good & Limited \\
LLM/RAG compatibility & Poor & Moderate & \textbf{Excellent} \\
Requires proprietary software & Partial & Yes & \textbf{No} \\
Long-term archival stability & High & Moderate & \textbf{High} \\
\bottomrule
\end{tabularx}
\end{table}

\section{Dataset and Ground Truth Preparation}

The test corpus consists of bachelor and master thesis documents. This scope
was chosen because academic theses are structurally demanding: they combine
dense body text, mathematical formulas, multi-column layouts, citations, tables,
and figures, often within the same document. They also represent the primary
use case this system targets. Testing across multiple thesis documents rather
than a single one allows RQ5 to check whether strategy rankings are stable
across documents or vary by the specific content mix of each thesis.

Documents were drawn from real sources rather than synthetic benchmarks.
Meuschke et al.\ \autocite{meuschke2023} found that tool performance varies
significantly by document structure, which is why the corpus was built to
include theses with different content mixes and layouts. The documents are
primarily in German, with a subset in English, reflecting the intended usage
context. Each selected document contains at least one structurally challenging
page: a formula-heavy section, a multi-column spread, or a table with complex
cell structure. The final corpus consists of three thesis documents totalling 90 evaluated
pages, 30 per document.

Each test document has a corresponding reference Markdown file, created by
hand. Each reference captures the full document structure: heading hierarchy,
paragraph flow, table cells, list nesting, and code blocks where present.
Pfitzmann et al.\ \autocite{pfitzmann2022} argue that human-created annotation is
the only reliable baseline for evaluating document parsing tools, because
automated references embed the same systematic errors as the tools being
tested. Creating the reference by hand ensures that the evaluation measures
real output quality rather than agreement between two imperfect automated
systems.

\section{Evaluation Metrics}
\label{sec:eval_metrics}

RQ4 asks how Markdown quality can be measured beyond BLEU and ROUGE.
Section~3.4 showed that flat sequence metrics reward correct content in
the wrong structure. The answer is to score each element type separately.
Seven metrics are used, grouped into three categories: content accuracy,
structural accuracy, and formatting accuracy.

Two metrics measure how much text was preserved. Text similarity uses
difflib's SequenceMatcher, which compares candidate and reference character
by character and returns a ratio between 0 and 1. Word overlap computes
the Jaccard similarity on the word sets of both outputs. Text similarity
catches character-level noise such as OCR artefacts. Word overlap catches
word-level omissions such as missing paragraphs.

Three metrics target document structure. Heading structure score compares
the candidate and reference heading counts at each level from H1 to H6,
skipping levels that are absent in both. List structure score checks that
bullet and numbered lists appear at the correct nesting level. Paragraph
count score compares paragraph counts with a 20 per cent tolerance: a
candidate within that range of the reference scores 1.0, with a linear
penalty beyond it. A document can have high text similarity and still fail
all three of these if the extractor collapses the heading hierarchy into
plain text.

Two metrics cover formatting elements. Table detection score checks for the
presence and count of pipe tables. Code block score checks for fenced code
blocks. Both are low-frequency elements, but a single missed table is a
meaningful failure for any document that contains one. These metrics make
that failure visible rather than diluted by high text similarity on the
surrounding prose.

TEDS (Tree Edit Distance based Similarity) offers finer-grained table
evaluation by comparing cell-level tree structure \autocite{zhong2020}. It was
considered but not used. The test corpus does not include tables complex
enough to require cell-level scoring, and producing cell-level annotations
by hand for a multi-document corpus carries a high annotation cost.
Table detection score covers the relevant case for this evaluation.

\section{Experimental Design}

The experiment applies each of the four strategies to every document in
the corpus at a fixed temperature of 0.0. Deterministic output is
required here because the evaluation compares strategies on the same
pages: any randomness introduced by higher temperatures would obscure
whether differences in score come from the strategy or from sampling
noise.

One vision-language model is used across all VLM strategies:
Qwen2.5-VL~7B, running locally through LM Studio. The text strategy
does not use a model at all. It calls PyMuPDF directly through the
pymupdf4llm library, which converts the PDF text layer to Markdown
without invoking any LLM. All four strategies therefore run under
identical hardware and model conditions, which keeps the comparison fair.

Text-only extraction at temperature 0.0 is the baseline. It is the
simplest possible extraction method: raw text from PyMuPDF with no
image input and no LLM call. A strategy must outperform it on at least
two metric categories to justify its additional computational cost.
Processing time is recorded per page, providing the raw measurements
for RQ3. Table~\ref{tab:strategy_comparison}
summarises the expected strengths and weaknesses of each approach based on
the literature \autocite{meuschke2023, whitington2012}.

Results are reported as mean scores per document category across all pages
in that category. Meuschke et al.\ \autocite{meuschke2023} used the same
category-level reporting and found that tools which rank well on one
document type often rank poorly on another. Category-level means make this
pattern visible. A strategy that wins overall but fails on specific documents is a weaker
result than one that performs consistently across the full corpus.

\begin{table}[ht]
\centering
\caption{Preliminary comparison of PDF extraction strategies}
\label{tab:strategy_comparison}
\renewcommand{\arraystretch}{1.4}
\begin{tabularx}{\textwidth}{lXXXX}
\toprule
\textbf{Criterion} & \textbf{Text-only} & \textbf{Image-only} & \textbf{Hybrid} & \textbf{Adaptive} \\
\midrule
Works on scanned PDFs        & No      & Yes               & Yes          & Yes      \\
Preserves reading order      & Partial & Yes (model-dep.)  & Yes          & Yes      \\
Handles complex tables       & Poor    & Moderate          & Good         & Good     \\
Handles mathematical formulas & Poor   & Good              & Good         & Good     \\
Computational cost           & Low     & High              & \mbox{Medium-High}  & \mbox{Low-Medium} \\
Structural fidelity          & Low     & Moderate          & High         & High     \\
Dependency on LLM            & None    & Full              & Partial      & Selective \\
\bottomrule
\end{tabularx}
\end{table}
